% IEEE conference manuscript. Every result is a macro or table generated from results/ by
%   python tasks.py paper-tables   (-> paper/generated/numbers.tex, paper/generated/tables.tex)
% Dataset counts typed in the text are copied from results/data_stats.json and docs/phases/PHASE_1.md.
\documentclass[conference]{IEEEtran}
\usepackage[utf8]{inputenc}
\usepackage{amsmath,amssymb}
\usepackage{booktabs}
\usepackage{graphicx}
\usepackage{xcolor}
\usepackage{url}
\usepackage{cite}
% Generated by python tasks.py paper-tables from results/. Do not edit.
\providecommand{\pending}[1]{\textcolor{red}{[pending: #1]}}
\newcommand{\cvMacroFOne}{0.998 $\pm$ 0.004}
\newcommand{\testAcc}{1.000}
\newcommand{\extMacroFOne}{0.119}
\newcommand{\extMacroFOneCI}{[0.083, 0.155]}
\newcommand{\extFlagged}{96.9\%}
\newcommand{\testFlagged}{7.9\%}
\newcommand{\aurocNear}{0.954}
\newcommand{\aurocFar}{0.963}
\newcommand{\pestTestMap}{0.520}
\newcommand{\pestModel}{yolo11n}
\newcommand{\sevLesionIoU}{\pending{severity/validation.json}}
\newcommand{\delAucGradcam}{0.606}
\newcommand{\delAucCentre}{0.629}
\newcommand{\delAucRandom}{0.286}
\newcommand{\insAucGradcam}{0.821}
\newcommand{\insAucCentre}{0.807}
\newcommand{\insAucRandom}{0.285}
\newcommand{\ohiMaxTierChange}{12.1\%}
\newcommand{\ohiUniformTierChange}{42.3\%}
\newcommand{\fcPrevHighDays}{216.2}
\newcommand{\apiDualMs}{\pending{api/latency.json}}


\begin{document}

\title{Near-Perfect Benchmark Accuracy Is Not Field Readiness:\\
An Audited Decision-Support Pipeline for Mango Leaf Disease, Pests and Orchard Health}

\author{\IEEEauthorblockN{[Authors withheld for review]}
\IEEEauthorblockA{[Affiliation]}}

\maketitle

\begin{abstract}
We present an open, fully reproducible decision-support pipeline for mango orchards. It combines a leaf-disease
classifier with Grad-CAM explanations, a colour-based lesion-area estimator, a pest detector, a transparent
rule-based Orchard Health Index (OHI), weather rules taken from a national IPM bulletin, and treatment advice
verified against official sources.

Every reported number is produced by code in the repository. With leakage-controlled group splits on MangoLeafBD,
three ImageNet-pretrained backbones reach a cross-validated macro-F1 of \cvMacroFOne{} and are statistically
indistinguishable. On an independent external set, however, the selected model reaches a macro-F1 of only
\extMacroFOne{} \extMacroFOneCI{}. A calibrated low-confidence rule routes \extFlagged{} of those images to a
``retake photo'' path.

The pest detector reaches a test mAP$_{50:95}$ of \pestTestMap{}. Every class comes from a single non-mango source,
so its detections are labelled unverified and kept out of the OHI. A sensitivity analysis shows that OHI rankings
are robust to the weights, but its tier labels are not: up to \ohiMaxTierChange{} of scenarios change tier under
$\pm$20\% changes to a single weight.

We report these limitations as primary results and release the pipeline, the audit trail and the protocols for
the missing expert validations.
\end{abstract}

\begin{IEEEkeywords}
mango, plant disease, domain shift, calibration, explainability, decision support, reproducibility
\end{IEEEkeywords}

\section{Introduction}
Image-based plant-disease classifiers routinely report accuracies close to 100\% on public datasets
\cite{mohanty2016plantvillage,ferentinos2018deep}. Such figures are often taken as evidence that a tool is ready for
growers. Two failure modes undermine that inference:
\begin{itemize}
\item \emph{leakage} between training and test images, for example multiple photographs of the same leaf
\cite{kapoor2023leakage};
\item \emph{domain shift} between curated dataset photographs and field conditions.
\end{itemize}
Decision-support systems add a third risk: scores and treatment advice that look authoritative but were never
validated.

We built a complete mango-orchard pipeline and audited every component with the same rule: a claim is made only if
a file under \texttt{results/} supports it. Our contributions are:
\begin{enumerate}
\item a leakage-controlled benchmark of three backbones on MangoLeafBD \cite{ahmed2023mangoleafbd}, including group
cross-validation, an independent external set, calibration and out-of-distribution (OOD) analysis
(Section~\ref{sec:disease});
\item an OHI that excludes unassessed inputs instead of scoring them as healthy, together with a global sensitivity
analysis of its expert-assigned constants (Section~\ref{sec:ohi});
\item a back-test of published disease-favourable weather conditions on ten years of reanalysis data, with an
explicit list of what it can and cannot claim (Section~\ref{sec:forecast});
\item a treatment database checked against official extension sources and the Indian banned-pesticide list, and a
reproducible artefact in which all numbers, tables and references are generated and checked by code.
\end{enumerate}

\section{Related Work}
Deep CNNs for leaf-disease recognition were popularised on PlantVillage \cite{mohanty2016plantvillage}, and a large
architecture comparison followed \cite{ferentinos2018deep}. For mango, MangoLeafBD provides eight classes of leaf
images \cite{ahmed2023mangoleafbd}.

We compare EfficientNet \cite{tan2019efficientnet}, ResNet \cite{he2016resnet} and MobileNetV3
\cite{howard2019mobilenetv3} classifiers, and YOLO-family detectors \cite{redmon2016yolo,jocher2023ultralytics}
pretrained on COCO \cite{lin2014coco}.

For the evaluation, we use:
\begin{itemize}
\item temperature scaling \cite{guo2017calibration} and maximum-softmax OOD scoring \cite{hendrycks2017baseline};
\item Grad-CAM \cite{selvaraju2017gradcam}, evaluated with deletion/insertion curves \cite{petsiuk2018rise}, whose
pitfalls are discussed in \cite{hooker2019roar}, and with the pointing game \cite{zhang2018excitation};
\item McNemar's test \cite{dietterich1998tests} with Holm correction \cite{holm1979}, and sensitivity analysis
following \cite{saltelli2008gsa}.
\end{itemize}

\section{Data}
\textbf{Disease.} We use five MangoLeafBD classes: Anthracnose, Bacterial Canker, Powdery Mildew, Sooty Mold and
Healthy. Leaf Spot is excluded because no licensed training data exist for it. The authors note that the 4{,}000
images contain only about 1{,}800 distinct leaves.

We grouped copies of the same leaf in two steps:
\begin{enumerate}
\item DINOv2 retrieval \cite{oquab2024dinov2};
\item verification with one-to-one SIFT matching \cite{lowe2004sift} and a RANSAC similarity transform
\cite{fischler1981ransac}.
\end{enumerate}
All splits are group-aware. This leaves 2{,}486 images from 1{,}085 leaves, with a held-out test split of 239 images.

\textbf{External set.} 381 independent images of four disease classes form the external set. 972 candidate images
were excluded because they derive from MangoLeafBD.

\textbf{Pests.} Three public sets (CC BY 4.0) give 2{,}559 images in 1{,}400 groups. A host-plant review found:
\begin{itemize}
\item fruit flies photographed on traps;
\item mealybugs on grape and in stock photos;
\item scale insects on coffee leaves.
\end{itemize}
None shows mango, so all three classes are treated as \emph{proxy} data.

\section{Methods}
\subsection{Disease classifier}\label{sec:disease}
\textbf{Protocol.} Each backbone is trained with 5-fold group cross-validation. Fold $k$ is held out, fold $k{+}1$
is used for early stopping, and the selection rule (highest mean macro-F1) was fixed before training.

\textbf{Training.} 224\,px inputs, strong augmentation, AdamW, label smoothing, class weights, seed 42.

\textbf{Final models.} Each final model is trained on the train split and evaluated once on the test split, the
external set, a near-OOD set (three MangoLeafBD classes unknown to the model) and a far-OOD set (pest images). All
evaluation runs in fp32: fp16 autocast produced non-finite logits on OOD inputs.

\textbf{Temperature and threshold.} A temperature is fitted on the validation split. Because validation accuracy is
100\%, the fit is not identifiable: the NLL decreases monotonically as $T\to0$. We therefore deploy $T{=}1$; this
rule was added post hoc and is disclosed as such. The low-confidence threshold is the maximum-softmax value that
keeps 95\% of validation images. Images below it trigger a ``retake photo'' message, and their disease result
enters neither the OHI nor the treatment advice.

\subsection{Lesion area and explanations}
\textbf{Lesion area.} A rule-based HSV segmenter defines the leaf as the largest filled green component, grown
through attached lesion tissue. Lesions are brown pixels plus pixels that are dark relative to the leaf, which
captures the low-saturation sooty-mould coating.

\textbf{Explanations.} Grad-CAM \cite{selvaraju2017gradcam} on the last convolutional layer is scored with
deletion/insertion AUC \cite{petsiuk2018rise} against two baselines: random ranking and a centred Gaussian prior.

\subsection{Pest detector}
YOLOv8n and YOLO11n \cite{jocher2023ultralytics} start from COCO weights pinned by SHA-256. Each is trained once
(seed 42, deterministic mode) and selected on validation mAP$_{50:95}$. The deployment metadata marks every class
\texttt{host\_verified: false}.

\subsection{Orchard Health Index}\label{sec:ohi}
$\mathrm{OHI}=\mathrm{clip}(100-\sum_{i\in A}\tilde w_i X_i)$, where $X_i\in[0,100]$ are penalties for disease,
pest, weather and nutrient status, and $\tilde w_i=w_i/\sum_{j\in A}w_j$ renormalises the nominal weights
$(0.40,0.35,0.15,0.10)$ over the set $A$ of components actually assessed.

A component is excluded when:
\begin{itemize}
\item it has no model;
\item the prediction is low-confidence;
\item weather is offline;
\item no nutrient sensor exists;
\item no pest class is host-verified.
\end{itemize}
Without an image-based component, no OHI is reported. All constants are expert-assigned. We perturb them:
\begin{itemize}
\item one at a time, by $\pm$10/20\%;
\item jointly, with Dirichlet-distributed weights;
\item by component ablation.
\end{itemize}
The 239 scenarios combine test-split predictions, segmented lesion area, randomly paired pest annotations and
randomly drawn hours of weather.

\subsection{Weather rules and back-test}\label{sec:forecast}
Thresholds are quoted from the National Horticulture Mission IPM bulletin \cite{nhm2012ipmmango}:
\begin{itemize}
\item anthracnose: RH${\ge}$95\% at 25--30\,$^\circ$C for 12\,h;
\item powdery mildew: T$_{\min}$ 10--13\,$^\circ$C, T$_{\max}$ 27--31\,$^\circ$C and RH 82--91\%;
\item bacterial canker: RH${\ge}$90\% and 25--30\,$^\circ$C.
\end{itemize}
They are applied to hourly NASA POWER/MERRA-2 \cite{gelaro2017merra2} data for Krishnagiri, Tamil Nadu, for
2015--2024. We also compare them with the rules the prototype used before the audit.

\subsection{Treatment advice}
Every recommendation links to a TNAU page \cite{tnau2026mango} or to the NHM bulletin, quotes it verbatim, and is
checked against the CIB\&RC list of banned and restricted pesticides \cite{cibrc2026banned}.

\section{Results}
% Generated by python tasks.py paper-tables from results/. Do not edit.
\begin{table}[t]
\centering
\caption{Five-fold group cross-validation on MangoLeafBD (mean $\pm$ SD over folds). $^\dagger$selected by mean macro-F1. Latency: CPU, batch 1.}
\label{tab:cv}
\footnotesize
\setlength{\tabcolsep}{3pt}
\begin{tabular}{lcccc}
\toprule
Model & Accuracy & Macro-F1 & Params (M) & ms \\
\midrule
EfficientNet-B3$^\dagger$ & 0.999 $\pm$ 0.002 & 0.998 $\pm$ 0.004 & 10.7 & 48.7 \\
ResNet-50 & 0.999 $\pm$ 0.003 & 0.997 $\pm$ 0.007 & 23.5 & 48.5 \\
MobileNetV3-L & 0.998 $\pm$ 0.002 & 0.996 $\pm$ 0.004 & 4.2 & 19.7 \\
\bottomrule
\end{tabular}
\end{table}

\begin{table}[t]
\centering
\caption{Final models: test split ($n{=}239$), independent external set ($n{=}381$, 4 classes), calibration (ECE) and OOD detection (AUROC, max softmax). Flag: share of external images below the low-confidence threshold.}
\label{tab:final}
\footnotesize
\setlength{\tabcolsep}{3pt}
\begin{tabular}{lccccc}
\toprule
Model & Test acc. & External macro-F1 [95\% CI] & ECE test/ext. & AUROC near/far & Flag \\
\midrule
EfficientNet-B3 & 1.000 & 0.119 [0.083, 0.155] & 0.086 / 0.416 & 0.954 / 0.963 & 0.97 \\
ResNet-50 & 1.000 & 0.079 [0.056, 0.102] & 0.081 / 0.703 & 0.976 / 0.995 & 0.54 \\
MobileNetV3-L & 0.992 & 0.167 [0.118, 0.214] & 0.088 / 0.443 & 0.949 / 0.939 & 0.87 \\
\bottomrule
\end{tabular}
\end{table}

\begin{table}[t]
\centering
\caption{Single-run ablations of the EfficientNet-B3 recipe on CV fold 0 ($n{=}279$); McNemar exact test against the baseline.}
\label{tab:ablation}
\footnotesize
\setlength{\tabcolsep}{3pt}
\begin{tabular}{lccc}
\toprule
Variant & Macro-F1 [95\% CI] & $\Delta$ & $p$ \\
\midrule
baseline & 1.000 [1.000, 1.000] & -- & -- \\
basic augmentation & 0.980 [0.954, 1.000] & -0.020 & 0.25 \\
img size 256 & 0.992 [0.971, 1.000] & -0.008 & 1 \\
no pretraining & 0.649 [0.584, 0.701] & -0.351 & 3.5e-18 \\
\bottomrule
\end{tabular}
\end{table}

\begin{table}[t]
\centering
\caption{Pest detectors on the proxy dataset (test $n{=}379$). Per class: AP50-95 for Fruit Fly / Mealy Bug / Scale Insect. Each class comes from a single non-mango source.}
\label{tab:pest}
\footnotesize
\setlength{\tabcolsep}{3pt}
\begin{tabular}{lccccc}
\toprule
Model & Params (M) & Val mAP$_{50:95}$ & mAP$_{50}$ & mAP$_{50:95}$ & Per class \\
\midrule
YOLOv8n & 3.0 & 0.533 & 0.743 & 0.519 & 0.539 / 0.118 / 0.899 \\
YOLO11n & 2.6 & 0.537 & 0.744 & 0.520 & 0.542 / 0.121 / 0.896 \\
\bottomrule
\end{tabular}
\end{table}

\begin{table}[t]
\centering
\caption{Grad-CAM faithfulness (test split, $n{=}239$): mean AUC; last column: share of images where Grad-CAM beats the centre prior (Wilcoxon $p$).}
\label{tab:xai}
\footnotesize
\setlength{\tabcolsep}{3pt}
\begin{tabular}{lcccc}
\toprule
Curve & Grad-CAM & Random & Centre & Better than centre \\
\midrule
deletion & 0.606 & 0.286 & 0.629 & 0.59 (0.0007) \\
insertion & 0.821 & 0.285 & 0.807 & 0.62 (0.0001) \\
\bottomrule
\end{tabular}
\end{table}

\begin{table}[t]
\centering
\caption{OHI sensitivity ($n{=}239$ scenarios): Spearman $\rho$ with the nominal scores (worst case for one-at-a-time) and share of scenarios whose tier changes.}
\label{tab:ohi}
\footnotesize
\setlength{\tabcolsep}{3pt}
\begin{tabular}{lcc}
\toprule
Perturbation & $\rho$ & Tier changed \\
\midrule
One weight $\times$ 0.8--1.2 & 0.994 & 0.121 \\
Dirichlet moderate (c=100) (median) & 0.996 & 0.075 \\
Dirichlet uniform Dirichlet(1,1,1) (median) & 0.831 & 0.423 \\
Dirichlet wide (c=20) (median) & 0.983 & 0.172 \\
Without disease & 0.759 & 0.481 \\
Without pest & 0.599 & 0.564 \\
Without weather & 0.987 & 0.188 \\
\bottomrule
\end{tabular}
\end{table}


\textbf{In-distribution benchmark.} All backbones are near-ceiling and statistically indistinguishable under
cross-validation (Table~\ref{tab:cv}). The only ablation with a measurable effect is removing ImageNet pretraining
(Table~\ref{tab:ablation}).

\textbf{Domain shift.} The test-split accuracy of \testAcc{} does not transfer (Table~\ref{tab:final}): external
macro-F1 is \extMacroFOne{} \extMacroFOneCI{} for the selected model. The low-confidence rule flags \testFlagged{} of
test images but \extFlagged{} of external images. In practice, the deployed system therefore asks for a retake
rather than issuing a confident wrong diagnosis on most out-of-domain photos. OOD AUROC is \aurocNear{} (near) and
\aurocFar{} (far).

\textbf{Pests.} The two detectors are equivalent (Table~\ref{tab:pest}). Fruit Fly and Scale Insect score high, but
each is a single homogeneous source, so background cues are confounded with the class. Mealy Bug, the heterogeneous
source, is poor.

\textbf{Explanations.} Grad-CAM beats the centre prior on both curves, but only modestly, on about 60\% of images
(Table~\ref{tab:xai}). Random-order deletion collapses the prediction faster than any saliency order. We attribute
this to the out-of-distribution artefact described in \cite{hooker2019roar}, so we do not use random deletion as a
lower bound.

\textbf{OHI.} Rankings are robust to the weights: $\rho\ge0.99$ for $\pm$20\% changes to a single weight
(Table~\ref{tab:ohi}). Tier labels are not: up to \ohiMaxTierChange{} of scenarios change tier one weight at a time,
and \ohiUniformTierChange{} (median) under uninformative weights.

\textbf{Weather rules.} In ten years of hourly reanalysis data, the bulletin's ``High'' conditions never occur. RH
${\ge}95\%$ is common, but never between 25 and 30\,$^\circ$C. The pre-audit rules raised a High alert on
\fcPrevHighDays{} days per year.

\section{Discussion and Limitations}
\textbf{Field readiness.} The central finding is negative. A leakage-controlled, near-perfect MangoLeafBD
benchmark says little about field photographs. Calibrated abstention mitigates the risk of confident errors, but it
does not create accuracy; field data from Indian orchards are required.

\textbf{Pest detector.} It cannot be claimed to work on mango trees.

\textbf{OHI and weather rules.} Their tiers and alerts rest on expert constants and on thresholds stated for other
regions. Coarse reanalysis data cannot resolve canopy leaf wetness, and no outbreak records were available.

\textbf{Pending validations.} Expert lesion masks (lesion area and the pointing game) and expert OHI ratings are
pending. Protocols and code are released with the artefact.

\textbf{Reduced scope.} ConvNeXt and Swin backbones, multi-seed ablations and a third detector were planned but not
run on the available 4\,GB GPU. We report no results for them.

\section{Conclusion}
Auditing every claim against generated evidence changed the story of this system: from ``accurate and ready'' to
``accurate in-distribution, safely abstaining out of distribution, and waiting for field validation''. We release
the pipeline, all results and the validation protocols so that the missing evidence can be collected.

\bibliographystyle{IEEEtran}
\bibliography{refs}

\end{document}
